Scientific applications execute on heterogeneous systems with multiple accelerators and deep memory and storage hierarchies, where data may traverse GPU memory, CPU memory, node-local storage, burst buffers, and parallel file systems. As data moves across these tiers of memory and storage, it is increasingly common for applications to desire in-flight processing applied along the data path, including compression, filtering, format conversion, encryption, and feature extraction for downstream analysis~\cite{ecp2019eqsim, mccallen2021eqsim, mccallen2024regional, patwary2015bdcats, tan2022analysis}. Applying these computations while data is in motion can reduce storage pressure, lower data movement cost, improve security, enable in-transit analysis, and prepare data for subsequent analysis without requiring a separate post-processing stage. Efficient and transparent in-transit processing is therefore becoming an important requirement for modern HPC data management systems.

Existing in-situ and in-transit systems have made significant progress toward moving computation closer to data. Existing in-situ systems such as DataSpaces~\cite{dataspaces}, FlexPath~\cite{flexpath}, and ADIOS~\cite{godoy2020adios2} support tight coupling, staging, and derived data computation alongside running simulations. In-transit approaches such as Runway~\cite{ravi2023runway} and the derived data work of Gainaru et al.~\cite{gainaru2024toderive} further explore runtime-level decisions about when and where to execute computations during data movement. 
However, these systems still contain important limitations. We define  ``in-flight processing'' as computations 
invoked implicitly 
by the applications or workflows, represented as isolated processing steps, or limited to a single transformation or analysis type. %As a result, 
With the current in-transit and in-situ processing, developers remain responsible for orchestrating multi-stage pipelines, tracking intermediate data states, and coordinating execution across heterogeneous resources. Runway \cite{ravi2023runway} is a recent example of an in-transit processing system, but it focuses on scheduling individual transformations rather than providing a composable pipeline abstraction. 
%We discuss these limitations in more detail in Section 2.

{\color{red}Here, we define \textbf{transformation} as an operation that mutates an object's state in place: the object keeps one name and identity throughout, only its on-storage representation changes (e.g., uncompressed and compressed, or unencrypted and encrypted, representations of the same object). We define \textbf{analysis}, by contrast, as an operation that derives new values from old values: it consumes one or more existing objects and produces one or more genuinely new, independently named objects with their own identity. Computing curl from a velocity field, for instance, yields a curl (i.e., the tendency of a vector field to spin or rotate around a specific point) object, not another representation of the velocity data. The two compose: once analysis produces a new object, that object is ordinary like any other and can itself be transformed, for example, by compressing a derived curl-magnitude object.}

Building a general in-flight processing system requires addressing several design challenges. First, the correct transformation depends on the current state of the data. For example, data that has already been compressed must be handled differently from data that remains uncompressed, and different portions of a dataset may exist in different states during a workflow. {\color{red}Analysis introduces a related problem: the system must track which derived objects have already been computed from which inputs, and whether a given derived result should be persisted or discarded once consumed.} The system therefore needs a way to track data states and identify valid transformation {\color{red}and analysis} paths. Second, modern HPC platforms include heterogeneous execution resources, and the best device for a given computation may depend on runtime conditions such as device utilization, transformation {\color{red}or analysis} cost, and data movement overhead. Static assignment of transformations {\color{red}or analysis operations} to CPUs or GPUs cannot adapt to this variability. Third, transformations {\color{red}and analysis operations} must be integrated with asynchronous data movement without stalling forward progress. Naively applying computations synchronously during I/O would disable the benefits of asynchronous execution. {\color{red}Fourth}, the system must shield application developers from this complexity entirely, requiring that transformation {\color{red}and analysis} pipelines be specified declaratively and executed transparently, without explicit orchestration on every data transfer. {\color{red}Finally, analysis operations often require combining multiple independently produced inputs before they can execute, and may in turn produce multiple derived outputs from a single computation. The system must therefore resolve when all required inputs for a given analysis operation are available, before triggering execution, and must route each resulting output to its own destination without requiring the application to track these dependencies itself.}

Our work targets modern, data-intensive scientific workflows where both computation and I/O are large and where in-flight processing costs are non-trivial~\cite{ecp2019eqsim, mccallen2021eqsim, mccallen2024regional, patwary2015bdcats, tan2022analysis}. For example, Vector Particle-In-Cell (VPIC)~\cite{tan2022analysis} is a plasma simulation code where the number of particles ranges from one trillion to ten trillion, and each simulation step can involve thousands of seconds of computation, leaving ample opportunity to hide checkpointing costs behind ongoing computation. Another representative example is the DOE EQSIM earthquake simulation workflow~\cite{ecp2019eqsim, mccallen2021eqsim, mccallen2024regional}, which couples extreme-scale computation with the need to efficiently transfer, store, and analyze large volumes of simulation output.

We implement our in-flight processing strategy, Data Flyway, in the Proactive Data Containers (PDC)~\cite{pdc, pdcdocs} object storage system, leveraging its client and server runtime, object/region abstractions, and storage-tier-aware data movement capabilities.

\paragraph{Contributions.} The contributions of this paper are as follows.
\begin{itemize}
\item We introduce Data Flyway, a declarative, data-centric transformation  {\color{red}and analysis} model in which applications specify the desired data state rather than explicitly invoking transformations {\color{red}or analysis operations}, addressing the need for developer transparency and decoupling transformation {\color{red}and analysis} intent from execution.
\item We propose a graph-based representation of transformation {\color{red}and analysis} pipelines, where transformations {\color{red}and analysis operations} are modeled as state transitions in a directed graph attached to I/O resources, enabling data state tracking and flexible, composable multi-stage data processing.
\item We design a runtime system for adaptive in-flight processing scheduling that dynamically determines when and where to execute computation based on resource utilization and workload characteristics, addressing heterogeneous device scheduling and enabling effective overlap with asynchronous data movement.
\item We implement our proposed Data Flyway framework in PDC and demonstrate that declarative in-flight processing graphs can be integrated into a production data management system to enable transparent, in-transit data processing, while noting that the design is generalizable to other object-based data management systems.
\item {\color{red}We introduce a mechanism for composing multiple independently produced inputs into a single analysis operation and routing its resulting outputs to their respective destinations, addressing the dependency resolution required by multi-input, multi-output analysis pipelines.}
\end{itemize}

The remainder of this paper is organized as follows: In Section 2 we discuss the background and motivation of this study, including related work and the PDC implementation context. In Section 3, we present the key components of Data Flyway. We present the extensive experimental results in Section 4. Finally, Section 5 outlines future directions and conclusions.